\label{int:chap:producto-nativo}
\index{produit natif}

L'extension euclidienne résidu--quotient d'APP contient deux applications
locales : l'application additive normalise les collisions et la multiplicative
produit des produits partiels. Dans cette section, on construit la multiplication
globale des mots en utilisant exclusivement ces deux applications. La correspondance avec la multiplication entière sera une
conséquence de l'entrelaceur, et non la définition de l'opération.

\section{Coordonnées euclidiennes de la cellule locale}

\begin{definition}[Applications additive et multiplicative sur un même support]
Soit $\Dnine=\{1,\ldots,9\}$. Pour $a,b\in\Dnine$, on définit par division
euclidienne avec reste dans $\Dnine$ :
\begin{align}
 a+b&=9q^+(a,b)+R^+(a,b),\label{int:eq:local-plus}\\
 ab&=9q^\times(a,b)+R^\times(a,b).\label{int:eq:local-times}
\end{align}
En particulier, le reste nul est représenté par $9$ et le quotient est réduit
d'une unité. Le quadruplet
\[
 \CellRQ(a,b)=
 (R^+(a,b),q^+(a,b);R^\times(a,b),q^\times(a,b))
\]
est l'extension arithmétique locale par résidu et quotient.
\end{definition}

L'agrégation selon les trois classes modulo trois du tableau visible
\(R^\times\) définit une matrice distincte des applications locales :
\[
(M_\Pi)_{rs}
=\frac19
\sum_{\substack{a\equiv r\;(3)\\b\equiv s\;(3)}}R^\times(a,b),
\qquad r,s\in\{1,2,0\},
\]
et, dans cet ordre des classes,
\[
M_\Pi=
\begin{pmatrix}4&5&6\\5&4&6\\6&6&9\end{pmatrix}.
\]
Ainsi, \(q^+\) et \(q^\times\) sont des coordonnées de quotient d'une cellule,
tandis que \(M_\Pi\) est un opérateur agrégé sur les classes résiduelles.

Les $81$ cellules satisfont exactement
\[
\begin{array}{c|rrrr}
 &\sum R&\sum q&\sum(9q+R)\\\hline
 +&405&45&810\\
 \times&459&174&2025
\end{array}
\]
et donc
\[
 2025-810=(459-405)+9(174-45)=54+9\cdot129=1215.
\]
Le quotient $q^\times$ est nécessaire à la reconstruction. Par exemple,
\[
 2\cdot5=9\cdot1+1.
\]
Si l'on ne conserve que $R^\times$, le produit $10$ se confond avec la valeur
visible $1$.

\subsection{Cocycles d'associativité et de distributivité}

La cohérence du transducteur peut déjà être décidée dans la cellule. Pour tout
$a,b,c\in\Dnine$, les trois identités sont vérifiées
\begin{align}
q^+(a,b)+q^+(R^+(a,b),c)
&=q^+(b,c)+q^+(a,R^+(b,c)),\label{int:eq:cocycle-plus}\\
q^\times(R^\times(a,b),c)+cq^\times(a,b)
&=q^\times(a,R^\times(b,c))+aq^\times(b,c),
\label{int:eq:cocycle-times}\\
q^\times(a,R^+(b,c))+aq^+(b,c)
&=q^+(R^\times(a,b),R^\times(a,c))\notag\\
&\quad+q^\times(a,b)+q^\times(a,c).
\label{int:eq:cocycle-distributive}
\end{align}

\begin{theorem}[Cohérence locale de l'extension euclidienne]
Les identités \eqref{int:eq:cocycle-plus}--\eqref{int:eq:cocycle-distributive},
avec les égalités correspondantes de résidus, expriment
respectivement l'associativité additive, l'associativité multiplicative et la
distributivité dans l'extension résidu--quotient. Les trois familles présentent
zéro écart sur les $9^3=729$ triplets.
\end{theorem}

\begin{proof}
En développant $(a+b)+c$ au moyen de \eqref{int:eq:local-plus}, le coefficient de
neuf est le membre de gauche de \eqref{int:eq:cocycle-plus} ; en développant
$a+(b+c)$, on obtient celui de droite. L'unicité de la division avec reste dans
$\Dnine$ égalise aussi bien les résidus que les quotients. Le même argument appliqué à
$(ab)c=a(bc)$ donne \eqref{int:eq:cocycle-times}. Enfin, les deux développements de
$a(b+c)=ab+ac$ produisent \eqref{int:eq:cocycle-distributive}. L'énumération
directe des \(729\) cas de chaque identité confirme le calcul
algébrique.
\end{proof}

\section{Bijectivité de la numération}

Le symbole $9$ n'est pas identifié à zéro. Dans $\CellRQ$, $9$ représente un
franchissement complet de la frontière nonadique. Le remplacement par les chiffres
$0,\ldots,8$ éliminerait le quotient euclidien qui intervient dans
l'entrelacement.

\begin{definition}[Mots canoniques]
Soit
\[
 \Wnine^+=\bigsqcup_{\ell\geq1}\Dnine^\ell,
 \qquad \Wnine=\{\varnothing\}\sqcup\Wnine^+.
\]
Les mots s'écrivent à partir du chiffre le moins significatif. Nous définissons
\[
 \operatorname{ev}_9(u_0,\ldots,u_{\ell-1})
 =\sum_{j=0}^{\ell-1}u_j9^j,
 \qquad \operatorname{ev}_9(\varnothing)=0.
\]
\end{definition}

\begin{theorem}[Unicité de la forme normale]
L'évaluation $\operatorname{ev}_9$ est une bijection de $\Wnine^+$ sur
$\N_{\geq1}$, et de $\Wnine$ sur $\N_0$.
\end{theorem}

\begin{proof}
Les mots de longueur $\ell$ couvrent l'intervalle
\[
 \frac{9^\ell-1}{8}
 \leq \operatorname{ev}_9(u)
 \leq 9\frac{9^\ell-1}{8}.
\]
Le minimum de l'intervalle suivant est supérieur d'une unité au maximum du
précédent. À l'intérieur de chaque intervalle, la division de $n-1$ par neuf détermine
de manière unique le chiffre
\[
 d=(n-1)\bmod9+1
\]
et le quotient qui est codé récursivement. Le processus se termine parce que le
quotient décroît strictement.
\end{proof}

Exemples :
\[
 9\longleftrightarrow(9),\qquad
 10\longleftrightarrow(1,1),\qquad
 81\longleftrightarrow(9,8),\qquad
 1000\longleftrightarrow(1,3,3,1).
\]

\section{Fibres des projections résiduelles dans les relèvements entiers}

Pour ce bloc, nous distinguons deux alphabets. La forme normale
\(\Wnine\) utilise les chiffres \(\{1,\ldots,9\}\), écrits à partir du moins
significatif. En revanche, \(\mathscr B_9\) désigne les écritures positionnelles
conventionnelles en base neuf, avec les chiffres \(\{0,\ldots,8\}\) et le plus
significatif à gauche. Les deux types ne sont pas identifiés.

Soit
\[
 \mathscr W_{10}^{(2)}=\{0,\ldots,9\}^2,\qquad
 \rho_{10}(a,b)=(b,a),\qquad
 \operatorname{ev}_{10}(a,b)=10a+b,
\]
et soit \(q_9:\mathbb N\to\mathbb Z/9\mathbb Z\) la réduction résiduelle.

\begin{lemma}[Descente résiduelle du renversement décimal]
\label{int:pf84:SYN30-RHO-DESC}
Le renversement sur les mots décimaux de longueur deux est involutif et son
ombre modulo neuf est l'identité :
\[
 \rho_{10}^2=I,\qquad
 q_9\!\left(\operatorname{ev}_{10}(\rho_{10}(a,b))\right)
 =
 q_9\!\left(\operatorname{ev}_{10}(a,b)\right).
\]
La longueur fixe conserve les zéros initiaux.
\end{lemma}

\begin{proof}
La double permutation retrouve \((a,b)\). Comme \(10\equiv1\pmod9\),
\[
 10b+a\equiv a+b\equiv10a+b\pmod9.
\]
Un mot dont le double renversement ne serait pas lui-même ou dont la somme des chiffres
changerait modulo neuf réfuterait le lemme.
\end{proof}

\begin{lemma}[Descente résiduelle du carré]
\label{int:pf84:SYN30-SQ-DESC}
L'application \(s(n)=n^2\) induit l'endoapplication résiduelle
\[
 [n]_9\longmapsto[n]_9^2,
 \qquad q_9(s(n))=q_9(n)^2.
\]
\end{lemma}

\begin{proof}
La compatibilité d'une congruence avec le produit donne
\(n\equiv m\pmod9\Rightarrow n^2\equiv m^2\pmod9\). Un entier pour lequel
\(q_9(n^2)\ne q_9(n)^2\) réfuterait l'énoncé.
\end{proof}

\begin{proposition}[Fibres résiduelles du renversement entier]
\label{int:pf84:SYN30-RHO-NONFAITH}
L'ombre résiduelle ne reconstruit pas le relèvement entier
\(\widehat\rho_{10}=\operatorname{ev}_{10}\circ\rho_{10}\). En effet,
\[
 18\equiv27\pmod9,\qquad
 \widehat\rho_{10}(1,8)=81\ne72=\widehat\rho_{10}(2,7).
\]
\end{proposition}

\begin{proof}
Les deux entiers appartiennent à la classe résiduelle zéro, mais leurs renversements sont
distincts. Il n'existe donc pas de
\(\overline\rho:\mathbb Z/9\mathbb Z\to\mathbb N\) qui factorise
simultanément ces deux valeurs.
\end{proof}

\begin{proposition}[Fibres résiduelles du carré entier]
\label{int:pf84:SYN30-SQ-NONFAITH}
La même ombre ne reconstruit pas non plus le carré entier :
\[
 18\equiv27\pmod9,\qquad
 18^2=324\ne729=27^2.
\]
\end{proposition}

\begin{proof}
Une fonction \(\overline s:\mathbb Z/9\mathbb Z\to\mathbb N\) qui
factoriserait \(s\) devrait attribuer à la même classe zéro les deux valeurs
\(324\) et \(729\), ce qui est impossible.
\end{proof}

\begin{corollary}[Témoin croisé \(27\)--\(72\)--\(729\)]
\label{int:pf84:SYN30-CROSS-WITNESS}
\[
 \widehat\rho_{10}(2,7)=72,\qquad
 27^2=729=9^3=3^6,
\]
et, dans l'alphabet conventionnel \(\mathscr B_9\),
\[
 27=30_9,\qquad72=80_9,\qquad729=1000_9.
\]
Les trois entiers ont une ombre résiduelle zéro et une racine numérique positive neuf.
\end{corollary}

\begin{proof}
Les deux opérations se calculent directement ; les divisions successives par
neuf produisent les trois écritures. Le résultat échouerait si l'on confondait
\(\mathscr B_9\) avec \(\Wnine\), si l'une des écritures était incorrecte
ou si l'on affirmait que la descente résiduelle fidèle reconstruit l'opération et sa
généalogie complète.
\end{proof}

\section{Additionneur natif et propagation du quotient}

\begin{algorithm}[Somme de mots]
On parcourt deux mots par positions, avec une retenue initiale \(c_0=0\). Si, à
une position, il y a deux chiffres, on applique \eqref{int:eq:local-plus} ; s'il n'y en a qu'un,
ce chiffre est pris comme résidu provisoire et le quotient provisoire est zéro.
Lorsque \(c_j\ne0\), on applique une seconde cellule additive au résidu provisoire
et à \(c_j\). On publie le nouveau résidu et l'on transporte la somme des
quotients à la position suivante. S'il reste à la fin une retenue non nulle, elle est
publiée comme dernier chiffre. De cette manière, les cellules locales ne reçoivent jamais le
symbole absent \(0\), qui n'appartient pas à \(\Dnine\).
\end{algorithm}

\begin{lemma}[Invariant de l'additionneur]
Si $u\boxplus_\#v$ est le mot produit par l'algorithme précédent,
\[
 \operatorname{ev}_9(u\boxplus_\#v)
 =\operatorname{ev}_9(u)+\operatorname{ev}_9(v).
\]
La retenue intermédiaire est dans $\{0,1,2\}$.
\end{lemma}

\begin{proof}
Après traitement des positions $0,\ldots,k-1$, soit $c_k$ le quotient en attente.
Les identités locales impliquent
\[
 \sum_{j<k}(u_j+v_j)9^j
 =\sum_{j<k}w_j9^j+c_k9^k,
\]
en incluant comme zéro les chiffres absents. Chaque nouvelle cellule conserve cette
égalité. Le maximum local est $9+9+2=20$, de sorte que $c_{k+1}\leq2$. À la
fin, le dernier quotient est publié, et l'identité globale demeure.
\end{proof}

\section{Multiplication par un chiffre}

\begin{algorithm}[Transducteur scalaire]
$d\in\Dnine$ étant fixé, chaque chiffre $u_j$ produit
\[
 u_jd=9q_j+r_j,
 \qquad
 (r_j,q_j)=\bigl(R^\times(u_j,d),q^\times(u_j,d)\bigr).
\]
Soit $c_j\in\{0,\ldots,9\}$ le quotient entrant. Comme $0\notin\Dnine$,
la transition est définie par branches :
\[
 (w_j,c_{j+1})=
 \begin{cases}
 (r_j,q_j),&c_j=0,\\[1mm]
 \bigl(R^+(r_j,c_j),q_j+q^+(r_j,c_j)\bigr),&c_j\in\Dnine.
 \end{cases}
\]
Ainsi, aucune application locale n'est évaluée hors de son domaine. La première
branche exprime l'absence de quotient en attente ; la seconde normalise la
collision au moyen de l'application additive.
\end{algorithm}

\begin{lemma}[Invariant scalaire]
Le mot $d\odot_\#u$ produit par le transducteur satisfait
\[
 \operatorname{ev}_9(d\odot_\#u)
 =d\operatorname{ev}_9(u).
\]
Le quotient multiplicatif appartient à $\{0,\ldots,9\}$.
\end{lemma}

\begin{proof}
L'argument est celui de l'additionneur, en remplaçant chaque terme local par
$u_jd$. Si $c_j=0$, l'identité
$u_jd=9q_j+r_j$ conserve directement l'égalité partielle. Si
$c_j\in\Dnine$, la seconde division euclidienne donne
\[
 r_j+c_j=9q^+(r_j,c_j)+R^+(r_j,c_j),
\]
et conserve la même égalité après transport de
$q_j+q^+(r_j,c_j)$. Le maximum $9\cdot9+9=90$ et la forme du reste dans
$\Dnine$ maintiennent le quotient dans la borne déclarée. Le dernier quotient est
publié comme chiffre.
\end{proof}

\section{Produit global par Horner}

\begin{definition}[Produit natif]
Soit $v=(v_0,\ldots,v_{m-1})$. En partant de $a_m=\varnothing$, on définit pour
$j=m-1,\ldots,0$
\[
 a_j=(9\odot_\#a_{j+1})\boxplus_\#(v_j\odot_\#u).
\]
Le résultat est
\[
 u\boxtimes_\#v:=a_0.
\]
\end{definition}

Cette définition ne contient pas
\[
 \operatorname{dec}_9(\operatorname{ev}_9(u)\operatorname{ev}_9(v)).
\]
L'évaluation n'apparaîtra que dans la démonstration de l'entrelacement.

\begin{theorem}[Multiplicativité native]
Pour tous les mots $u,v\in\Wnine$,
\begin{equation}
 \label{int:eq:ev-multiplicative}
 \boxed{
 \operatorname{ev}_9(u\boxtimes_\#v)
 =\operatorname{ev}_9(u)\operatorname{ev}_9(v).}
\end{equation}
Le mot vide est absorbant et $(1)$ est l'unité.
\end{theorem}

\begin{proof}
Par les deux lemmes précédents,
\[
 \operatorname{ev}_9(a_j)
 =9\operatorname{ev}_9(a_{j+1})
  +v_j\operatorname{ev}_9(u).
\]
L'itération de Horner donne
\[
 \operatorname{ev}_9(a_0)
 =\operatorname{ev}_9(u)
  \sum_{j=0}^{m-1}v_j9^j.
\]
Le dernier facteur est $\operatorname{ev}_9(v)$. L'unicité de la forme normale
identifie le mot résultant.
\end{proof}

\begin{corollary}[Monoïde multiplicatif de mots]
La structure
\[
\WordMonoid=(\Wnine,\boxtimes_\#,(1))
\]
est un monoïde commutatif ; le mot vide est son élément absorbant. Sa
construction n'utilise pas la multiplication globale des entiers.
\end{corollary}

L'exemple de frontière le plus court est
\[
 (9)\boxtimes_\#(9)=(9,8),
 \qquad 9+8\cdot9=81.
\]
Les \(17\) cellules locales de la multiplication produisent
\[
 (9,8)\boxtimes_\#(9,8)=(9,8,8,8),
 \qquad \operatorname{ev}_9(9,8,8,8)=6561.
\]

\section{Linéarisation matricielle de l'état généalogique et construction de l'entrelaceur complètement positif}

\begin{definition}[Secteurs de Hilbert]
Soient
\[
 \Krad=\ell^2(\Wnine^+),
 \qquad
 \Hplus=\ell^2(\N_{\geq1}).
\]
Nous définissons sur les bases canoniques
\[
 W_\times e_u=e_{\operatorname{ev}_9(u)}.
\]
Dans \(\Hplus\), l'opérateur nombre est noté
\[
 \mathsf N e_n=n e_n,
 \qquad \Spec(\mathsf N)=\N_{\geq1}.
\]
Le vide peut être ajouté séparément comme vecteur de vide $e_0$ ; il ne se
confond pas avec $\Hplus$, dont le spectre commence à un.
\end{definition}

\begin{theorem}[Unitarité et carré multiplicatif]
$W_\times$ est unitaire. Dans le cœur algébrique
$\mathbb C[\Wnine^+\times\Wnine^+]$, soit
\[
 \mu_\#(e_u\otimes e_v)=e_{u\boxtimes_\#v},
 \qquad
 \mu_{\N}(e_a\otimes e_b)=e_{ab}.
\]
Alors on a le carré
\[
\begin{CD}
\Krad\otimes_{\mathrm{alg}}\Krad @>{W_\times\otimes W_\times}>>
\Hplus\otimes_{\mathrm{alg}}\Hplus\\
@V{\mu_\#}VV @VV{\mu_{\N}}V\\
\Krad @>>{W_\times}> \Hplus
\end{CD}
\]
et, par conséquent, le carré commute.
\end{theorem}

\begin{proof}
La bijection des formes normales envoie une base orthonormale sur une autre et prouve
l'unitarité. Sur un tenseur de base, les deux chemins produisent
$e_{\operatorname{ev}_9(u)\operatorname{ev}_9(v)}$ par
\eqref{int:eq:ev-multiplicative}. La linéarité donne l'égalité dans le cœur.
\end{proof}

Les applications $\mu_\#$ et $\mu_{\N}$, définies initialement sur les cœurs
algébriques respectifs, sont densément définies et non bornées. Pour
\(c=(c_{u,v})\in\ell^2(\Wnine^+\times\Wnine^+)\), chaque somme sur une fibre de produit
est finie. L'adjoint de \(\mu_\#\) contient tous les vecteurs à support
fini dans son domaine, de sorte que \(\mu_\#\) est fermable. Le domaine de sa
fermeture est
\[
 \Dom(\overline{\mu_\#})=
 \left\{c\in\ell^2(\Wnine^+\times\Wnine^+):
 \sum_{w\in\Wnine^+}
 \left|\sum_{u\boxtimes_\#v=w}c_{u,v}\right|^2<\infty
 \right\}.
\]
Chaque fibre est finie parce que $\operatorname{ev}_9$ l'identifie aux couples
ordonnés de diviseurs d'un entier. La même description vaut pour
\(\overline{\mu_{\N}}\) après application de \(W_\times\otimes W_\times\), et
$W_\times$ entrelace les graphes des deux fermetures.

\section{Coproduit et sélecteur natifs}

L'ouverture multiplicative est maintenant définie avant de regarder le codomaine :
\[
 \Delta_\#e_w
 =\sum_{u\boxtimes_\#v=w}e_u\otimes e_v.
\]
La somme s'obtient en énumérant les produits du transducteur. On ne demande pas si
un entier ordinaire en divise un autre.
Dans le secteur entier, nous définissons, sur le cœur algébrique correspondant,
\[
 \Delta_\times e_n
 =\sum_{\substack{a,b\in\N_{\geq1}\\ab=n}}e_a\otimes e_b.
\]
Les deux coproduits énumèrent des couples ordonnés.

\begin{theorem}[Coassociativité native]
Dans le cœur algébrique,
\[
 (\Delta_\#\otimes I)\Delta_\#
 =(I\otimes\Delta_\#)\Delta_\#.
\]
De plus,
\[
 (W_\times\otimes W_\times)\Delta_\#
 =\Delta_\times W_\times.
\]
\end{theorem}

\begin{proof}
Les deux membres de la première égalité énumèrent une fois chaque triplet natif
$a\boxtimes_\#b\boxtimes_\#c=w$. La seconde égalité découle de la bijection
entre fibres induite par \eqref{int:eq:ev-multiplicative}.
\end{proof}

\begin{theorem}[Coproduit arithmétique complet]
\label{int:pf84:C06}
Dans le cœur algébrique
\(c_{00}(\mathbb N_{\geq1})\subset\ell^2(\mathbb N_{\geq1})\), soit
\[
 \Delta_\times e_n=\sum_{ab=n}e_a\otimes e_b.
\]
Sa fermeture a pour domaine
\[
 \operatorname{Dom}(\overline{\Delta_\times})
 =
 \left\{x=(x_n):
 \sum_{n\geq1}d(n)|x_n|^2<\infty\right\},
\]
est coassociative dans le cœur et, pour \(Qe_n=ne_n\), satisfait
\[
 (\log Q\otimes I+I\otimes\log Q)\Delta_\times
 =\Delta_\times\log Q.
\]
\end{theorem}

\begin{proof}
Les supports de \(\Delta_\times e_n\) sont orthogonaux pour des entiers distincts
et leur norme au carré est le nombre \(d(n)\) de diviseurs, ce qui donne le domaine
pondéré de la fermeture. Les deux coassociations réindexent la même somme
sur \(abc=n\). Enfin, \(\log a+\log b=\log n\) dans chaque terme
prouve l'entrelacement. Une triple factorisation comptée de manière différente
ou l'échec de cette identité réfuterait le théorème.
\end{proof}

\begin{theorem}[Coproduit réduit et sélecteur d'irréductibles]
\label{int:pf84:C07}
Dans le secteur entier, le coproduit réduit
\[
 \Delta_\times^{\rm red}e_n
 =\sum_{\substack{ab=n\\a,b>1}}e_a\otimes e_b
\]
est fermable. L'opérateur
\[
 N_\times^{\rm red}
 =\bigl(\overline{\Delta_\times^{\rm red}}\bigr)^*
  \overline{\Delta_\times^{\rm red}}
\]
est positif autoadjoint et diagonal :
\[
 N_\times^{\rm red}e_n=d_{\rm red}(n)e_n.
\]
Son noyau est engendré par \(e_1\) et par les \(e_p\) avec \(p\) irréductible ;
par conséquent,
\[
 P_{\Irr}
 =\mathbf1_{\{0\}}(N_\times^{\rm red})
  -\lvert e_1\rangle\langle e_1\rvert
\]
sélectionne les irréductibles et exclut l'unité.

Dans le secteur natif, le coproduit correspondant est
\[
 \widetilde\Delta_\#e_w
 =\sum_{\substack{u\boxtimes_\#v=w\\u,v\ne(1)}}e_u\otimes e_v.
\]
L'entrelaceur \(W_\times\) transporte
\(\Delta_\times^{\rm red}\), \(N_\times^{\rm red}\) et \(P_{\Irr}\) vers
\(\widetilde\Delta_\#\),
\[
 A_\#:=
 \bigl(\overline{\widetilde\Delta_\#}\bigr)^*
 \overline{\widetilde\Delta_\#},
 \qquad
 P_{\Irr,\#}
 =(I-P_{(1)})\mathbf1_{\{0\}}(A_\#),
 \qquad
 P_{(1)}=\lvert e_{(1)}\rangle\langle e_{(1)}\rvert.
\]
\end{theorem}

\begin{proof}
Les supports des images de deux vecteurs de base distincts sont orthogonaux.
La norme au carré de
\(\Delta_\times^{\rm red}e_n\) est \(d_{\rm red}(n)\), qui s'annule
exactement pour \(n=1\) ou \(n\) irréductible. Cela prouve la diagonalisation
et le noyau déclaré ; soustraire le projecteur de \(e_1\) élimine l'unité.
La multiplicativité de \(W_\times\) identifie les fibres entières aux
fibres natives et transporte les trois opérateurs. Un composé dans le noyau,
un irréductible en dehors de celui-ci ou l'inclusion de l'unité dans \(P_{\Irr}\)
réfuteraient le théorème.
\end{proof}

\begin{proposition}[Coïncidence nonadique des nombres premiers et composés]
\label{int:pf84:B05}
Soit \(\gcd(n,90)=1\),
\(\tau_n=\operatorname{ord}_n(10)\) et
\[
 M_n=\frac{10^{\tau_n}-1}{n}.
\]
Alors \(M_n\equiv0\pmod9\), aussi bien pour les nombres premiers que pour les composés. En
particulier, la clôture nonadique et la période décimale prises isolément ne sont pas un
sélecteur d'irréductibles.
\end{proposition}

\begin{proof}
Comme \(9\mid10^{\tau_n}-1\) et \(n\) est inversible modulo neuf, la
divisibilité par neuf survit au quotient. Les nombres premiers \(7,13\) et les
composés \(77,91,143\) partagent la période six ; les nombres de Carmichael
apportent des contrôles supplémentaires. L'énoncé ne serait réfuté que si la
congruence ou la période séparaient tous les nombres premiers de tous les composés
dans le domaine déclaré.
\end{proof}

Cette affirmation est strictement arithmétique : elle identifie des fibres de
factorisation et les opérateurs qu'elles induisent sur le secteur des formes
normales.

\ifdefined\HMTInternalTraceability
\section{Extension monoïdale du produit natif aux histoires compatibles}
\label{int:pf84:SYN-05D-HIST}

Le résultat distingue l'espace \(\Krad\) des formes normales canoniques et
l'espace \(\mathcal K_{\#,\mathrm{hist}}\) des histoires complètes de
cellules et de quotients propagés.
De nombreuses histoires peuvent produire la même forme normale. Par conséquent,
$W_\times$ est unitaire sur $\Krad$, mais la projection d'une histoire sur
sa valeur empêche qu'il soit injectif sur toutes les classes de transport.

Soit \(\mathscr X_n^{\rm hist}\) l'ensemble des histoires enrichies de
profondeur \(n\) et soit
\(\mathcal H_n=\ell^2(\mathscr X_n^{\rm hist})\). Notons \(P_n\) la
projection orthogonale sur le sous-espace engendré par les histoires
survivantes et supposons définie une multiplication \(\mu_n\) sur un
domaine algébrique de
\(\mathcal H_n\otimes_{\mathrm{alg}}\mathcal H_n\). Une extension monoïdale
du produit de mots au système inverse consiste en une famille
d'applications \(\iota_n:\Krad\to\mathcal H_n\) qui commute avec les
restrictions de profondeur et satisfait, pour chaque \(n\),
\[
 P_n\,\iota_n\mu_\#
 =P_n\,\mu_n
 (\iota_n\otimes\iota_n)
\]
dans le cœur algébrique. Cette égalité est la condition de compatibilité qui
définit l'extension. Les théorèmes précédents n'établissent ni l'existence ni
l'unicité d'une famille \((\iota_n)_n\) ; ils démontrent le produit dans le
monoïde \(\WordMonoid\) et déterminent l'équation qu'une extension à l'espace
des histoires doit satisfaire. Le transfert conserve donc le
statut de \emph{programme ouvert} : il n'hérite pas de l'exactitude du produit dans
\(\WordMonoid\). L'inexistence d'une famille compatible, l'échec de la
commutation avec les troncatures, la perte de survie, de quotient,
de retenue, d'orientation ou de mémoire, ou le non-respect de l'identité précédente
sont ses falsificateurs.
\fi

\section{Énumération exhaustive et nécessité des coordonnées structurelles}

L'énumération finie comprend :
\[
\begin{array}{@{}lr@{}}
\toprule
\text{Test}&\text{Résultat}\\\midrule
\text{Cellules par application locale}&81\\
\text{Parcours réversibles de numérotation}&59050\\
\text{Paires de l’entrelaceur}&532900\\
\text{Erreurs de l’entrelaceur}&0\\
\text{Triplets associatifs}&531441\\
\text{Erreurs associatives}&0\\
\text{Fenêtre de coproduit}&1\text{ à }500\\
\text{Irréductibles obtenus}&95\\
\bottomrule
\end{array}
\]

\begin{definition}[Quotient multiplicatif]
Conserver seulement $R^\times$ transforme $2\cdot5=10$ en la valeur visible $1$.
Perturber $q^\times(9,9)$ de $8$ à $7$ transforme $81$ en $72$. Le quotient
local est nécessaire.
\end{definition}

\begin{definition}[Application additive]
L'application multiplicative produit des produits partiels, mais ne résout pas leurs collisions.
Éliminer $q^+$ rompt la normalisation globale ; par exemple, elle cesse de reproduire
les quotients intermédiaires de $10\cdot18=180$.
\end{definition}

\begin{definition}[Coordonnées et symbole de bord]
Le zéro ordinaire n'est pas un chiffre de la construction. Le mot vide est
$0$ et le symbole $9$ conserve sa valeur et son quotient. Mélanger les coordonnées
positionnelles indexées à partir de zéro avec les étiquettes $1,\ldots,9$ modifie les
cellules.
\end{definition}

\begin{proposition}[Portée exacte des ablations et du recensement fini]
\label{int:pf84:D10}
Les ablations précédentes prouvent, dans les conventions déclarées :
\begin{enumerate}
  \item sans \(q^\times\), \(2\cdot5=10\) se réduit au résidu visible \(1\) ;
  \item remplacer \(q^\times(9,9)=8\) par \(7\) transforme \(81\) en \(72\) ;
  \item sans \(q^+\), la normalisation des collisions additives échoue.
\end{enumerate}
L'énumération exhaustive comprend les \(2\cdot81=162\) cellules des deux
applications, les \(729\) triplets locaux, \(532\,900\) paires de mots et
les \(95\) irréductibles inférieurs ou égaux à \(500\). Dans ces domaines finis,
les identités indiquées sont démontrées, mais la survie
coinductive de toute histoire TPK n'est pas prouvée.
\end{proposition}

\begin{proof}
Les trois contre-exemples se calculent directement avec
\eqref{int:eq:local-plus} et \eqref{int:eq:local-times}. L'énumération des quatre
domaines finis produit les cardinaux écrits dans l'énoncé et n'ajoute
aucune identité coinductive à celles qui ont déjà été démontrées dans le monoïde
de mots.
\end{proof}

\section{Structure algébrique du monoïde de mots}

La construction démontre, dans le monoïde \(\WordMonoid\) :
\begin{enumerate}
 \item une forme normale unique pour chaque entier non négatif ;
 \item une somme et un produit globaux construits uniquement avec les deux applications locales ;
 \item la conservation symbolique de la valeur sous la propagation du quotient ;
 \item un unitaire $W_\times$ et le carré multiplicatif exact ;
 \item un coproduit diviseur et un sélecteur irréductible définis nativement ;
 \item la nécessité de conserver la classe de transport lorsque l'on quitte le
 secteur des formes normales.
\end{enumerate}

Le passage à l'espace enrichi se réduit maintenant à un problème d'existence :
construire une famille \((\iota_n)_n\) compatible avec les restrictions de
profondeur et démontrer qu'elle satisfait l'équation de la section précédente.
Cette question est distincte de l'entrelaceur exact sur
\(\WordMonoid\), qui est déjà déterminé par la forme normale et les
applications locales.
